% !TeX root = ../papers/nc-p3.tex
\section{階数一曲線の総次数と曲線外の例外平面}
\label{sec:rank-one-curve-degree-bound}
\begin{theorem}\label{thm:rank-one-curve-degree-bound}
前節と同じ被約窓で、階数一軌跡の曲線部分の被約台を $C$ とする。その総次数は6以下である。
総次数6の場合は連結成分が三つ以上あり、曲線外の指定印は高々一つである。
特に既約な階数一曲線の次数は5以下であり、平面的なものは直線か二次曲線に限られる。
\end{theorem}
\begin{proof}[一般列の有限基点]
入力・出力の二つの二次Veroneseの同一の交叉を読むと $C$ には出力曲線 $C'$ との同型があり、
それぞれの超平面束 $L,L'$ が $L^2=(L')^2$ を満たす。特に各成分の次数は両側で等しい。
$Nv=0$ の全入射は次元3以下である。階数一は前定理で高々曲線でその上の繊維は平面、
階数二は高々曲面で繊維は直線、階数三は高々曲面で繊維は一点である。
従って一般の $v$ について四つの独立な二次式 $Nv$ の共通基点は有限である。

この有限基点スキームの長さは6以下である。一般の三つの線型結合が長さ8の二次式完全交叉を切る。
これを避ける無限遠平面を選んでアフィン化すると、その有限環の次数濾過の随伴環は三変数の二次完全交叉である。
実際三つの最高次式は無限遠で共通零点がなく、随伴環への全射が両辺の長さ8により同型になる。
従って最後の次数3の一次元商を読む汎関数は完全な積対合を与え、次数2以下を消す
（先頭次数の元には随伴環の完全対合で相補次数の元を選べる）。
第四の二次式の像は非零である。完全交叉の飽和イデアルの次数二は最初の三つが張るからである。
もしその生成イデアルが長さ1なら全ての倍元はその定数倍になり全ての跡が消える。完全性に反するので
残余の基点長は $8-2$ 以下である。一般の出力平面 $v^\perp$ は $C'$ を $d=\deg C$ 個の異なる点で切り、
対応する入力点が全て基点になる。これで $d\le6$ を得る。
平面曲線なら二次Veroneseの包のベクトル次元は高々6であり、その出力曲線も平面的である。
非退化整曲線なら四つの一次座標の積は少なくとも七次元だからである（一点での相異なる位数の和を取る）。
出力平面の零化列 $v$ に対する $Nv$ はその曲線上で消えるので、
定理\ref{thm:opposite-column-divisor}の固定列の曲線次数上界2が適用される。
\end{proof}
\begin{proof}[重複因子は二種類共存できない]
行列式に相異なる二つの重複既約因子があると仮定する。次数4以上の重複因子は、
四次の平方の場合に全二十印で偏導三次式が零となるため不可能である。
次数3も残余二次式外の六印以上で特異となって不可能である。
これらは定理\ref{thm:low-rank-surface-factors}の最初の段落の議論であり階数一の有限性を使わない。
従って各重複因子は一次か二次である。
二次二つなら式は二既約二次式の積の平方となる。両者の共通点以外の印は一方の頂点に限られ、
少なくとも十八の印が一つの二次式に載ってしまう。
一次と二次なら残余二次を含めた全印での重複度の総和は $2\cdot7+2\cdot15+15=59$ 以下で、必要な60に足りない。
一次二つならその両平面の外に六印以上あり、残余四次式はそこで三重になる。
既約四次なら三重印は共線で高々三つ、二既約二次の積なら高々二頂点であるから、残余四次には平面因子がある。
残った三次の重複度の和は高々24（可約でも一次7、二次15の上界を使う）。従って全体で
$2\cdot7+2\cdot7+7+24=59$ 以下になりこれも矛盾する。
特に階数二以下の曲面成分は存在してもただ一つである。
\end{proof}
\begin{proof}[一般列の例外となる平面]
$d=6$ とする。任意の曲線外の階数一の入力点 $p$ に対応する出力点を $y$ とすれば $y\notin C'$ である。
$y$ を通る一般平面 $v^\perp$ は $C'$ をなお六つの異なる点で切る。実際この平面系は $C'$ 上基点がなく、
$y$ からの射影は各曲線成分上非定数である（$y$ を通る直線成分はない）。
標数零なので一般切断は横断的である。従って $Nv=0$ はその六点と $p$ により七つ以上の解を持ち、
先の長さ上界により正次元になる。$v$ の動く平面 $\Pi_y=y^\perp$ を支配する入射成分で
その一般繊維が正次元のものを取る。これは元の全入射の三次元成分であり、出力像はちょうど $\Pi_y$。
階数一曲線上の成分ではない。その曲線の点 $p'$ ごとの核平面が全て $\Pi_y$ なら、対応する出力点が
一定となって閉埋込みに反するからである。従って階数二の曲面上の核直線束の閉包である。
一つの既約曲面はこの三次元成分を一つだけ持ち、その出力像も一意である。
先に曲面成分は高々一つと示したので、異なる $y$ を持つ曲線外の点は高々一つとなる。

$C$ の連結成分数を $b$ とする。$h^0(C,L L')\le 2d+b$ である。
各既約成分の正規化上では次数に1を加えたものが上界で、連結成分内で一本ずつ新しい曲線をつなぐ際に
共有点での値の一致が少なくとも一条件を与える。各枝上の切断は共有点で生成するためである。
一方、入力・出力の十六の双一次式は全印上で独立である。
$y^tBx=0$ が全印で成立すれば $N(x)Bx$ は三次評価の単射性で恒等零、一般可逆性から $B=0$ だからである。
曲線外の印は高々一つなので $16\le 2d+b+1=13+b$。よって $b\ge3$ である。
\end{proof}

\begin{corollary}\label{cor:degree-six-columns-double-cover}
総次数6の場合、一般の $v$ に対する有理写像 $x\mapsto[N(x)v]$ は次数二で三空間を支配する。
\end{corollary}
\begin{proof}
六つの異なる基点が全長を使い切り、各点で四式は極大イデアルを生成する。一般の三結合はそれらの一次部を張り
各基点で重複度1を与える。完全交叉の残る長さ二は一般の目標点の逆像である。
\end{proof}

\begin{proposition}\label{prop:degree-six-rank-one-component-partitions}
総次数6なら曲線上の印はちょうど十九、外は一つである。曲線の既約成分の次数分割は
$(4,1,1),(3,2,1),(3,1,1,1)$ のいずれかである。
四次成分があれば滑らかな有理曲線で十三印、三次成分があれば有理正規三次で十印を含む。
各直線には三印、各二次曲線には六印があり、異なる成分の印は重ならない。
\end{proposition}
\begin{proof}
まず任意の滑らかな平面二次曲線上の印は六個以下である。隣接する辺の六つの関係を
その曲線と隣の三空間との積 $\mathbb P^1\times\mathbb P^3$ へ制限すると $(2,1)$ の独立な六式を得る。
依存があれば元の非零関係が平面の式と一つの隣接一次式との積になり、有限辺の純関係の不在に反する。
共通零点は有限なので一般の四結合が長さ8の完全交叉を切る。この上で消える $(2,1)$ 式の空間は
その四結合だけである。Koszul分解を $(2,1)$ で捩ると二段目以降の第二因子は
$\mathcal O(-1),\mathcal O(-2),\mathcal O(-3)$ で全コホモロジーが消える。
残る二式の評価は従って独立で、それらの生成イデアルは長さ2以上を持つ。よって切断は長さ6以下である。

$d=6$ で連結成分が三つ以上なので各既約成分の次数は高々4である。
直線と二次曲線の印数上界はいずれもその次数の三倍である。
三次・四次の整曲線では三次評価から $3\deg+1$ が上界であり、三次以上の成分は高々一つである。
実際二つあれば他の成分に正次数を残せない。従って十九印以上に達するには三次か四次が一つ必要で、
全ての上界で等号となる。分割は記載の三つのみである。四次成分の場合、十三印での三次評価の全射性から
$h^0(C_4,\mathcal O(3))\ge13$ である。正規化上で次数12の十三切断は種数零を強制する
（種数が正なら一点の単純極だけを持つ非定数関数はなく、点を順に足せば上界は12）。
従って全十三切断が正規化から降り、その完全系が正規化を埋め込むので曲線自体が滑らかである。
三次成分は前定理で平面的でなく、最小次数の有理正規三次である。
\end{proof}
